\begin{table}[htbp]\centering\small
\caption{Pensiones: el perímetro adjudicado y sus vehículos, línea P (mdp)}\label{cua:pens}
\begin{tabular}{lrrr}
\toprule
Concepto & 2026 & 2027 & Var. real (\%) \\
\midrule
\textbf{Perímetro del Anexo 3, por diferencia} & \textbf{1.704.159,2} & \textbf{1.840.746,1} & \textbf{+3,9} \\
Pensiones y jubilaciones (clas. económica) & 1.704.200,0 & 1.840.700,0 & +3,9 \\
Ramo 19 Aportaciones a Seguridad Social & 1.541.518,7 & 1.661.007,9 & +3,6 \\
Ramo 20 Bienestar & 674.510,0 & 697.106,3 & $-$0,6 \\
\bottomrule
\end{tabular}
\\[2pt]\begin{minipage}{0.92\textwidth}\scriptsize Los dos primeros renglones son \textbf{el mismo objeto por dos caminos independientes} ---la diferencia de los dos renglones del Anexo 3 del decreto y la línea de la clasificación económica del CGPE--- y coinciden al décimo en los dos ejercicios. El Ramo 20 incluye la pensión no contributiva y \textbf{no es sumable} con los anteriores.\end{minipage}
\end{table}
